% ===== PRÉAMBULE CANONIQUE — à coller VERBATIM en tête de CHAQUE .tex =====
\documentclass[11pt,a4paper]{article}
\usepackage[utf8]{inputenc}\usepackage[T1]{fontenc}\usepackage{lmodern}
\usepackage[french]{babel}
\usepackage{amsmath,amssymb,mathtools}\usepackage{icomma}
\usepackage[version=4]{mhchem}          % \ce{...} : équations & formules
\usepackage{siunitx}\sisetup{locale=FR} % \qty{}{}, \si{}, \ang{}
\usepackage{chemfig}                    % \chemfig{...} : structures développées
\usepackage{xcolor}\usepackage{colortbl}
\usepackage{tikz}\usepackage{pgfplots}\pgfplotsset{compat=1.18}
\usepackage[most]{tcolorbox}\usepackage{enumitem}\usepackage{array,booktabs}
\usepackage[a4paper,margin=2.2cm]{geometry}
\usetikzlibrary{arrows.meta,calc,angles,quotes,positioning,patterns,backgrounds,shapes.geometric}
\definecolor{primary}{RGB}{222,49,99}\definecolor{primarydark}{RGB}{150,22,55}
\definecolor{defcol}{RGB}{0,114,178}\definecolor{methcol}{RGB}{52,92,148}
\definecolor{excol}{RGB}{0,135,90}\definecolor{attncol}{RGB}{200,40,55}
\tcbset{boxbase/.style={enhanced,breakable,boxrule=0.9pt,arc=2.5pt,
  left=3mm,right=3mm,top=2.3mm,bottom=2.3mm,fonttitle=\bfseries,coltitle=white}}
% --- boîtes TUTORIEL (noms EXACTS, ne pas renommer) ---
\newtcolorbox{cle}[1][]{boxbase,colback=defcol!6,colframe=defcol,
  title={\textsc{À retenir}\ifx\relax#1\relax\relax\else\textmd{ : \itshape #1}\fi}}
\newtcolorbox{methode}[1][]{boxbase,colback=methcol!7,colframe=methcol,sharp corners,
  title={\textsc{Méthode}\ifx\relax#1\relax\relax\else\textmd{ --- \itshape #1}\fi}}
\newtcolorbox{exemplebox}[1][]{boxbase,colback=excol!5,colframe=excol,
  title={\textsc{Exemple}\ifx\relax#1\relax\relax\else\textmd{ : \itshape #1}\fi}}
\newtcolorbox{piege}[1][]{boxbase,colback=attncol!6,colframe=attncol,
  title={\textsc{Piège}\ifx\relax#1\relax\relax\else\textmd{ : \itshape #1}\fi}}
\newtcolorbox{remarque}[1][]{boxbase,colback=black!4,colframe=black!45,
  title={\textsc{Remarque}\ifx\relax#1\relax\relax\else\textmd{ : \itshape #1}\fi}}
% --- boîtes EXERCICE / CORRIGÉ (noms EXACTS, auto-comptées) ---
\newtcolorbox[auto counter]{exo}[1][]{boxbase,colback=primary!4,colframe=primary,
  title={\textsc{Exercice}~\thetcbcounter\ifx\relax#1\relax\relax\else\textmd{ --- \itshape #1}\fi}}
\newtcolorbox[auto counter]{corrige}[1][]{boxbase,colback=excol!4,colframe=excol,
  title={\textsc{Corrigé}~\thetcbcounter\ifx\relax#1\relax\relax\else\textmd{ --- \itshape #1}\fi}}
\newcommand{\R}{\mathbb{R}}\newcommand{\N}{\mathbb{N}}\newcommand{\Z}{\mathbb{Z}}\newcommand{\ds}{\displaystyle}
% (un \usepackage{fancyhdr}+en-tête est possible ; il est ignoré côté web)
% =========================================================================
\begin{document}

\begin{corrige}[reconnaître la loi]
On repère la grandeur constante, ce qui fixe la relation.
\begin{enumerate}[label=\alph*)]
  \item $T$ constante $\Rightarrow$ \textbf{Boyle-Mariotte} : $p_1V_1=p_2V_2$.
  \item $p$ constante (piston libre) $\Rightarrow$ \textbf{Charles} : $\dfrac{V_1}{T_1}=\dfrac{V_2}{T_2}$.
  \item $V$ constant (récipient rigide) $\Rightarrow$ \textbf{isochore} :
        $\dfrac{p_1}{T_1}=\dfrac{p_2}{T_2}$.
  \item $p$ et $T$ constants $\Rightarrow$ \textbf{Avogadro} : $\dfrac{V_1}{n_1}=\dfrac{V_2}{n_2}$.
\end{enumerate}
\end{corrige}

\begin{corrige}[raisonnement proportionnel sans calcul détaillé]
\begin{enumerate}[label=\alph*)]
  \item $V$ constant : $p\propto T$. $T$ est divisée par $2$ $\Rightarrow$ $p$ est \textbf{divisée par $2$}.
  \item $p$ constante : $V\propto T$. $T$ doublée $\Rightarrow$ $V$ \textbf{doublé}.
  \item $T$ constante : $p\,V=$ cste (inverse). $p$ divisée par $2$ $\Rightarrow$ $V$ \textbf{doublé}.
  \item $p$ constante : $V\propto T$. $V$ doublé $\Rightarrow$ $T$ absolue \textbf{doublée}.
\end{enumerate}
\end{corrige}

\begin{corrige}[loi de Charles avec conversions]
$V_2=V_1\dfrac{T_2}{T_1}$ (ou $T_2=T_1\dfrac{V_2}{V_1}$), en kelvins.
\begin{enumerate}[label=\alph*)]
  \item $T_1=300$ K, $T_2=360$ K : $V_2=2{,}0\times\dfrac{360}{300}=\qty{2.4}{\litre}$.
  \item $T_1=400$ K, $T_2=300$ K : $V_2=4{,}0\times\dfrac{300}{400}=\qty{3.0}{\litre}$.
  \item $T_2=T_1\dfrac{V_2}{V_1}=250\times\dfrac{6{,}0}{5{,}0}=\qty{300}{\kelvin}$.
\end{enumerate}
\end{corrige}

\begin{corrige}[isochore et Avogadro]
\begin{enumerate}[label=\alph*)]
  \item $p_2=p_1\dfrac{T_2}{T_1}=1{,}0\times\dfrac{500}{250}=\qty{2.0}{atm}$.
  \item $T_2=T_1\dfrac{p_2}{p_1}=300\times\dfrac{3{,}0}{2{,}0}=\qty{450}{\kelvin}$.
  \item $V_2=V_1\dfrac{n_2}{n_1}=6{,}0\times\dfrac{0{,}40}{0{,}25}=\qty{9.6}{\litre}$.
\end{enumerate}
\end{corrige}

\end{document}
